\subsection{Hyperparameter Tuning}

Hyperparameter tuning is conducted to assess the sensitivity and
robustness of the proposed two-stage classifier, rather than to
aggressively maximize performance.

\subsubsection{Two-Stage Rule Parameters}
We evaluate the impact of the minimum rule purity $\tau_p$ and
minimum rule support $\tau_s$ on macro-averaged F1.
Figure~\ref{fig:purity_f1} and Figure~\ref{fig:support_f1} show that
performance increases monotonically with rule purity, while remaining
largely stable across a broad range of support thresholds.

\textbf{[INSERT FIGURE: Macro F1 vs Rule Purity]}
\textbf{[INSERT FIGURE: Macro F1 vs Rule Support]}

These results indicate a wide performance plateau, suggesting that
the rule-based stage is not overly sensitive to the exact choice of
thresholds once high-purity patterns are selected.

\paragraph{Linear Model Regularization}
The regularization parameter $C$ of the Stage~2 linear classifier is
varied across a broad range.
Figure~\ref{fig:c_f1} shows that macro-F1 remains stable for
$C \in [0.5, 1.5]$, with only marginal degradation for larger values.

\textbf{[INSERT FIGURE: Mean Macro F1 vs C]}

This behavior further confirms that the probabilistic stage operates
in a stable regime and does not require fine-grained tuning.

\paragraph{Best Configurations}
Table~\ref{tab:best_twostage} reports the top-performing configurations,
including the corresponding rule coverage.

\textbf{[INSERT TABLE: Top Two-Stage Configurations (C, purity, support, macro-F1, coverage)]}

Overall, the tuning results reveal a flat optimum across a wide
hyperparameter region.
This indicates that the observed performance gains are driven by the
two-stage formulation itself rather than by fragile parameter choices.
